\chapter{Surya Namaskara}
\label{ch:surya-namaskara}

\emph{What is the sun salutation, in what order are its positions done, and how does the breath move with them?}
Surya namaskara is a flowing sequence of postures joined with the breath. The candidate demonstrates it, names each position, and may be asked the mantra that goes with each. This chapter gives the twelve positions of the guidebook sequence, the breath and the mantras, the cautions, and the points where schools differ. The sequence follows the official guidebook; in the demonstration and the exam, use the Board's protocol sequence where it differs.

\section{What it is}

\emph{Surya namaskara} means salutation to the sun. It is counted as part of hatha yoga practice and is traditionally done in the early morning or the evening facing the sun, on an empty stomach; that is custom, not a rule. It is a sequence of asanas, not a single asana. Done slowly, it loosens the spine and the limbs and prepares the body for the asanas of the day, so it is usually done after the loosening practices and before the asanas.

It is traditionally credited with an energised and well-cleared body, flexibility and strength, better circulation and digestion, calm and focus, and help with weight control. These are well-being aims, not cures.

\section{The sequence}

The breath follows one simple rule: \emph{up and back, breathe in; down and forward, breathe out}. Each position has a mantra, said silently; the mantra is optional.

\begin{enumerate}
\item \textbf{Pranamasana} (prayer pose), \emph{Om Mitraya Namah}. Stand straight, palms together at the chest, and breathe normally.
\item \textbf{Hastottanasana} (arms raised), \emph{Om Ravaye Namah}. Breathe in, raise the arms and stretch back.
\item \textbf{Padahastasana} (standing forward bend), \emph{Om Suryaya Namah}. Breathe out, bend forward and place the palms by the feet. Bend the knees at first if needed.
\item \textbf{Ashwa sanchalanasana} (equestrian pose), \emph{Om Bhanave Namah}. Breathe in, step the right leg back, knee down, and look up.
\item \textbf{Parvatasana} (mountain, an inverted V), \emph{Om Khagaya Namah}. Breathe out, lift the hips, and take the heels towards the floor.
\item \textbf{Ashtanga namaskara} (eight-limb salute), \emph{Om Pushne Namah}. Breathe out and lower the body so that the two feet, two knees, two hands, the chest and the chin touch the floor, hips slightly raised.
\item \textbf{Bhujangasana} (cobra), \emph{Om Hiranyagarbhaya Namah}. Breathe in, raise the head and chest, and arch back.
\item \textbf{Parvatasana}, \emph{Om Marichaye Namah}. Breathe out and lift the hips.
\item \textbf{Ashwa sanchalanasana}, \emph{Om Adityaya Namah}. Breathe in, bring the right foot forward between the hands, and look up.
\item \textbf{Padahastasana}, \emph{Om Savitre Namah}. Breathe out and bend forward.
\item \textbf{Hastottanasana}, \emph{Om Arkaya Namah}. Breathe in, arms up and back.
\item \textbf{Pranamasana}, \emph{Om Bhaskaraya Namah}. Breathe out and return to the prayer pose.
\end{enumerate}

The sequence is a mirror: positions two to five come back in reverse order as positions eight to eleven. Only two positions are not repeated, ashtanga namaskara and bhujangasana, the sixth and seventh. One full round is the twelve positions done twice, the second time with the left leg going back. Rest afterwards in shavasana, relaxing each part from the toes to the head.

\ycbfig{ycb1-fig-suryanamaskar}
\ycbfig{ycb1-fig-sunsal1}
\ycbfig{ycb1-fig-sunsal2}

\section{How to practise}

Keep the movements slow and smooth, never jerky. Keep the breath steady and rhythmic, with no forceful breathing and no breath holding. Start with a few rounds and add slowly. Practise in a well-aired place.

\begin{caution}
Do not practise in acute high blood pressure, hernia or heart disease, or after a stroke; nor during fever or inflammation. Avoid it during menstruation. People with back problems practise only after advice from a yoga expert.
\end{caution}

\section{Where schools differ}

Several sequences are taught under the same name, and a candidate may meet them in other classes.
\begin{itemize}
\item \textbf{Names.} The guidebook's hastottanasana is the same as hasta uttanasana; its parvatasana is commonly called adho mukha svanasana; and ashtanga namaskara is also written sashtanga namaskara.
\item \textbf{Plank.} Some schools add a plank position (chaturanga dandasana) and so count three positions that are not repeated. The guidebook sequence has no plank. Do not mix the two in the exam.
\item \textbf{Mantra.} Some schools close with a single mantra; the guidebook gives the twelve names instead, each said silently with its position, and says that chanting is not compulsory.
\end{itemize}

\begin{recall}
\begin{itemize}
\item Salutation to the sun: a sequence of asanas joined with the breath, not a single asana.
\item Twelve positions, starting and ending in pranamasana; one round is two sets, one per leg.
\item Up and back, breathe in; down and forward, breathe out.
\item Not repeated: ashtanga namaskara and bhujangasana.
\item First mantra \emph{Om Mitraya Namah}, last \emph{Om Bhaskaraya Namah}.
\end{itemize}
\end{recall}

\subsection*{Worked examples}
\begin{itemize}
\item \textbf{A student learned a sequence with a plank at another centre.} Explain that both are taught, but the Board protocol sequence has no plank; in the exam, demonstrate the protocol sequence.
\item \textbf{A student with a stiff back cannot reach the floor in padahastasana.} Let them bend the knees and place the hands where they reach. Ask them to get advice from a yoga expert before doing full rounds.
\item \textbf{Model answer: which positions take an in-breath?} The ones that open the front of the body: hastottanasana, ashwa sanchalanasana and bhujangasana.
\end{itemize}

\subsection*{Check yourself}
\begin{enumerate}
\item What does surya namaskara mean?
\item Is surya namaskara a single asana?
\item How many positions are in one set of the guidebook sequence?
\item Name the first and last positions.
\item Which positions take an in-breath?
\item Which positions are not repeated in the second half?
\item What makes one full round?
\item What is the first mantra?
\item Give two cautions.
\end{enumerate}

\subsection*{Answers}
\begin{enumerate}
\item Salutation to the sun.
\item No, it is a sequence of asanas joined with the breath.
\item Twelve.
\item Pranamasana for both.
\item Hastottanasana, ashwa sanchalanasana and bhujangasana.
\item Ashtanga namaskara and bhujangasana.
\item The twelve positions done twice, once with each leg going back.
\item \emph{Om Mitraya Namah}.
\item Any two of: acute high blood pressure, hernia, heart disease, after a stroke, fever or inflammation, menstruation; back problems only after expert advice.
\end{enumerate}

\paragraph{Sources.} The official guidebook for Level I (2016, older scheme), chapter on surya namaskara; the Common Yoga Protocol (general practice guidelines); the Board syllabus for Level 1.
